\documentclass[11pt,twoside,openright]{book}

\usepackage[T1]{fontenc}
\usepackage[utf8]{inputenc}
\usepackage{amsmath,amssymb,amsthm,mathtools,bm}
\usepackage{mathpazo}
\usepackage{microtype}
\usepackage[inner=1.35in,outer=1.1in,top=1.25in,bottom=1.35in]{geometry}
\usepackage{booktabs}
\usepackage{fancyhdr}
\usepackage[hidelinks]{hyperref}
\usepackage{cleveref}

\setlength{\parindent}{1.4em}
\setlength{\parskip}{0pt}
\linespread{1.07}
\emergencystretch=3em

\pagestyle{fancy}
\fancyhf{}
\fancyhead[LE]{\small\thepage\quad\textsc{A Cinema of Individual Editions}}
\fancyhead[RO]{\small\nouppercase{\leftmark}\quad\thepage}
\renewcommand{\headrulewidth}{0.3pt}
\fancypagestyle{plain}{\fancyhf{}\fancyfoot[C]{\small\thepage}\renewcommand{\headrulewidth}{0pt}}

\theoremstyle{plain}
\newtheorem{theorem}{Theorem}[chapter]
\newtheorem{proposition}[theorem]{Proposition}
\newtheorem{lemma}[theorem]{Lemma}
\newtheorem{corollary}[theorem]{Corollary}
\theoremstyle{definition}
\newtheorem{definition}[theorem]{Definition}
\newtheorem{example}[theorem]{Example}
\theoremstyle{remark}
\newtheorem{remark}[theorem]{Remark}

\newcommand{\Capac}{\mathrm{Cap}}
\newcommand{\Rob}{\mathrm{Rob}}
\newcommand{\Inv}{\mathrm{Inv}}
\DeclareMathOperator{\loss}{loss}

\title{\Huge\bfseries A Cinema of Individual Editions\\[1.2ex]
\large Chapter Fifteen: Exclusion and Reference Corpora\\[0.4ex]
\large Working Draft}
\author{Flyxion\\[0.3ex]\small Independent Researcher}
\date{November 2026}

\begin{document}

\frontmatter
\maketitle
\tableofcontents
\mainmatter
\setcounter{chapter}{14}

\chapter{Exclusion and Reference Corpora}
\label{ch:exclusion}

A production may declare material that its editions must not reproduce. The declaration might protect a competitor's film, an earlier film of the same studio, or a performance that the studio has no right to repeat. The earlier chapters counted editions without regard to any such list. This chapter subtracts a list and asks what is left. It keeps the count apart from any determination of legal status. Whether an item is protected, whether two works are alike in the sense a court would apply, and whether a given edition infringes are questions that this book does not address and that no count in it answers.

The counts are taken on the lawn scene after the agreement rule, with the $238$ class tuples of production A and the $96$ of production B. An excluded item is a class tuple, and a reference corpus is a set of them. The corpora are drawn at random, either from the pool itself, which is the worst case for the subtraction, or from the whole product of classes, which is the case in which an outside corpus shares the vocabulary of the production but not its rules. Each count averages sixty draws.

\section{A Policy and Its Three Parts}\label{sec:policy15}

An exclusion policy has a corpus $X$, a representation $\phi$ in which an edition and an excluded item are compared, and a threshold $\rho$. An edition is \emph{excluded} if some item of the corpus lies within $\rho$ of it in the representation. The surviving editions are those outside every neighborhood.

The three parts are separate decisions, and the count depends on all of them. The corpus is a fact about the world. The representation is a choice about what counts as the same material, and the threshold is a choice about how close counts as too close. Exact subtraction is the case $\rho=0$ in the finest representation, in which an item is excluded only if it is identical.

\section{Exact Subtraction}\label{sec:exact15}

If the representation is the class tuple, or the take tuple, and the threshold is zero, each excluded item removes exactly itself. The count falls by the number of excluded items that were in the pool, and nothing else changes. Excluding a thousand take tuples of production A by digest leaves $56{,}834$ of $57{,}834$, and the same thousand leave $2{,}072$ of the $3{,}072$ of production B.

The class tuple is a coarser unit than the take tuple, and one class tuple stands for several take tuples. In A every class has three takes, so excluding one class tuple removes $3^5=243$ take tuples. In B it removes $2^5=32$. Excluding ten class tuples removes $2{,}430$ take tuples of A and $320$ of B. The difference between the two units is the number of takes per class raised to the number of slots, and which unit a policy names decides how much it takes.

\section{Approximate Subtraction}\label{sec:approx15}

An approximate policy removes a neighborhood around each item. In the Hamming representation on class tuples, the neighborhood of radius $\rho$ around an excluded tuple has, in the pool of A, a mean of $10.6$ tuples at radius one, with a range of eight to fourteen, and a mean of $47.4$ at radius two, with a range of thirty-two to seventy. In the pool of B it has exactly seven at radius one and exactly twenty-one at radius two. Radius two in a scene of five slots covers a fifth of the pool of A and a fifth of the pool of B around each excluded item.

\Cref{tab:surv15} gives the survivors when the corpus is drawn from the pool, for $n$ excluded items, under exact subtraction, under radius one and radius two, and under a coarser representation described below.

\begin{table}[ht]
\centering
\begin{tabular}{@{}llrrrr@{}}
\toprule
Production & Excluded items $n$ & Exact & $\rho=1$ & $\rho=2$ & Direction pattern\\
\midrule
A (238) & 1 & 237 & 227.2 & 188.9 & 228.9\\
 & 5 & 233 & 189.5 & 78.0 & 198.9\\
 & 10 & 228 & 149.2 & 29.4 & 160.8\\
 & 20 & 218 & 92.9 & 3.3 & 119.7\\
\midrule
B (96) & 1 & 95 & 89.0 & 75.0 & 48.0\\
 & 5 & 91 & 65.1 & 26.4 & 2.4\\
 & 10 & 86 & 43.8 & 7.0 & 0.0\\
 & 20 & 76 & 16.8 & 0.3 & 0.0\\
\bottomrule
\end{tabular}
\caption{Mean number of surviving class tuples after $n$ excluded items drawn from the pool, over sixty draws, by representation and threshold.}\label{tab:surv15}
\end{table}

Twenty excluded items leave production A with about three of its $238$ tuples at radius two, and leave production B with about seventeen of its $96$ at radius one. A modest corpus under a modest threshold exhausts a small production. The corpus drawn from the whole product of classes is less severe, because most of its items lie outside the pool: for $n=20$ the survivors of A are $235.7$, $196.3$ and $48.8$ at the three thresholds, and those of B are $95.6$, $88.9$ and $40.1$.

\section{The Representation Decides}\label{sec:repr15}

The last column of \Cref{tab:surv15} compares the editions by their direction pattern alone, which is the sequence of directions in which the character moves across the slots. Two editions are alike in this representation if they move the same way at every slot, whatever the footage. The pool of A has forty direction patterns, the largest shared by sixteen tuples. The pool of B has two, since B has no neutral class, and the direction rule forces every tuple to move all left or all right. Excluding one item of B in this representation excludes the forty-eight tuples that share its pattern, which is half the pool, and five items exclude all but $2.4$ on average.

The representation is a statement of what is taken to be the same material, and a coarse one takes more. The example does not suggest that the direction pattern is a good representation of anything a policy would want to protect. It shows that with five excluded items the survivors of B differ by a factor of about forty between the finest representation and one that sounds modest, and that the choice has to be made, and defended, before the count means anything.

\section{Exclusion and the Packing Number}\label{sec:pack15}

The packing numbers of Chapter~2 bound how many editions can be pairwise distinguishable. Removing a neighborhood from the pool lowers them, and the amount is bounded when the neighborhood is small relative to the tolerance.

\begin{proposition}[Packing under exclusion]\label{prop:packexcl}
Let $\mathcal P$ be a pool, let $\varepsilon\ge0$, and let $B(x,\rho)$ be a Hamming ball with $2\rho\le\varepsilon$. Then every set of editions that are pairwise more than $\varepsilon$ apart contains at most one edition of the ball, and hence the packing number of $\mathcal P\setminus B(x,\rho)$ at tolerance $\varepsilon$ is at least the packing number of $\mathcal P$ less one. Removing $n$ such balls lowers it by at most $n$.
\end{proposition}

\begin{proof}
Two editions of the ball are within $2\rho\le\varepsilon$ of each other by the triangle inequality, so a set of editions pairwise more than $\varepsilon$ apart contains at most one. A packing of $\mathcal P$ with its member in the ball removed is a packing of the remainder.
\end{proof}

The proof gives a more general bound. If the removed region can be covered by $c$ balls of radius at most $\varepsilon/2$, then at most $c$ points of any $\varepsilon$-packing lie in that region, and deleting them from a maximum packing leaves a packing of the survivors of size at least $M_\varepsilon-c$. The region may be a ball of any radius, since a larger ball can itself be covered by smaller ones. The difficulty is in getting a useful value of $c$ from the representation and the exclusion policy without enumerating the pool.

The proposition is a bound, and the lawn scene shows it is not tight. The packing number of A at tolerance two is sixteen, and after removing balls of radius one around one, three, five and ten excluded items it is still sixteen, against lower bounds of fifteen, thirteen, eleven and six. At tolerance one with exact subtraction it stays at fifty-four after ten exclusions. In B the packing number at tolerance two is twelve before exclusion, and twelve, eleven, ten and eight after one, three, five and ten, which is within the bound at each step, and at tolerance one with exact subtraction it stays at thirty-two. All of these were solved to optimality.

The reading is that a policy whose threshold is no more than half the tolerance of the promise costs the promise at most one distinguishable edition per excluded item, and in these examples much less. A policy with a larger threshold has no such bound, and the survivor counts of \Cref{tab:surv15} show how quickly it can take everything.

\section{Grammar Against Combination}\label{sec:grammar15}

Cinematic convention is reusable. A close shot of hands on a bolt, a wide view of a lawn, and a cut at a glance are not the property of any film, and a policy must be able to exclude a specific combination of imagery, performance, and timing without excluding the conventions from which every film is built. In the lawn scene a convention is a class at a slot, and a combination is a tuple.

Exact subtraction cannot damage a convention, since it removes tuples one at a time and the fewest tuples through any class in the pool of A is twenty-two, and sixteen in the pool of B. Approximate subtraction can, because a neighborhood removes every tuple through a class that its items cover. \Cref{tab:gram15} gives, for the pool of A and of B, the mean number of classes left with no surviving tuple and the number left with under five percent of their tuples. The pool of A has nineteen classes over the five slots and the pool of B twenty-two.

\begin{table}[ht]
\centering
\begin{tabular}{@{}llrrr@{}}
\toprule
Production & Threshold and $n$ & Classes with none left & Classes below 5\%\\
\midrule
A & $\rho=1$, $n=20$ & 0.00 & 0.00\\
 & $\rho=2$, $n=10$ & 1.53 & 5.80\\
 & $\rho=2$, $n=20$ & 10.48 & 17.57\\
\midrule
B & $\rho=1$, $n=20$ & 1.28 & 2.70\\
 & $\rho=2$, $n=10$ & 10.40 & 13.47\\
 & $\rho=2$, $n=20$ & 20.67 & 21.30\\
\bottomrule
\end{tabular}
\caption{Mean number of the classes of the pool that lose every tuple, and that fall below five percent of their tuples, after $n$ excluded items drawn from the pool, over sixty draws. Exact subtraction loses none in either production.}\label{tab:gram15}
\end{table}

At radius two, twenty excluded items leave production A with eight or nine of its nineteen classes still in use and leave production B with one of twenty-two. The policy that was meant to remove specific combinations has removed the grammar. A policy that must retain the conventions can state the requirement directly, as a minimum number of tuples that every class must keep, and check it after subtraction. The test is cheap, since it is a count through each class of the survivors, and a policy that fails it has a threshold that is too wide for the scene.

\section{The Count Apart From the Law}\label{sec:law15}

Every number in this chapter is a count of tuples under a declared policy. None of them is a determination that an excluded item is protected, that a neighborhood is a fair measure of likeness, or that an edition outside every neighborhood is free of any claim. The corpus is declared by someone, the representation is chosen by someone, and the threshold is set by someone, and the law's own standard of likeness does not have to coincide with any of the three. What the arithmetic supplies is the cost of a declaration: how much of a production a given policy takes, how few distinguishable editions it leaves, and whether it keeps the conventions from which films are made. That cost is the thing a production can know before deciding.

\section{What the Chapter Establishes and What It Assumes}\label{sec:limits15}

The proposition on packing holds for any pool and any Hamming ball. The counts rest on assumptions stated beside them: that the excluded items are class tuples of the same scene, that the Hamming dissimilarity on classes stands in for a measured one and a direction pattern for a coarse one, that the corpora are random draws and so show typical and not worst cases, and that a convention is a class at a slot, which is a convention of this toy and no more. A real corpus is not a set of tuples of the production's own vocabulary, and a real representation of imagery, performance, and timing is a feature extractor and a threshold that Chapter~2 posed and did not measure.

The next chapter turns from subtraction to generation. Chapter~16 asks what a generator must represent to produce takes under the same rules, and the rejection of samples in a generator is the exclusion of this chapter paid for in compute.

The script \texttt{verify\_ch15.py} reproduces each figure of this chapter.

\begin{thebibliography}{99}

\bibitem{flyxion2026gf}
Flyxion.
\newblock \emph{Generative Film: Categories, Sheaves, and Types for Cinema}.
\newblock Unpublished monograph, October 2026.

\end{thebibliography}

\end{document}
